%%
%%

\subsection{Entscheidungsregel}
Lernziele:

Entscheidungsregel um die $k$-Grenze und somit Annahme- und
Ablehnugsbereich festzulegen.
\newpage 

Wir gehen wieder von unserm Pharma-Unternehmen aus, das mit ihrem Produkt bei 80\,\% der Patienten Linderung verspricht.
Nun haben wir lediglich bei 63 Personen aus 100 eine Linderung festgestellt und wir \textbf{vermuten}, dass die 80\,\% Regel nicht stimmt.

Immer wenn wir etwas \textbf{vermuten}, \textbf{glauben} oder \textbf{behaupten}, dann haben wir sofort eine Alternativhypothese:


\begin{definition}{Alternativhypothese $H_1$}{}
Unter der \textbf{Alternativhypothese} verstehen wir eine Behauptung, welche die Nullhypothese in Frage stellt.
\end{definition}

Unsere Alternativhypothese könnte lauten:

\begin{center}Das Medikament liefert in weniger als 80\,\% Linderung.\end{center}

\begin{bemerkung}{Alternativhypothese}{}
  Die Alternativhypothese beinhaltet oft ein \textbf{weniger als} oder ein \textbf{mehr als}. In der Regel wird nicht gesagt, \textbf{wie viel} weniger oder mehr --- dies wird durch das Signifikanzniveau festgelet. Natrürlich kann auch ein \textbf{nicht} enthalten sein, doch bei useren einseitigen Hypothesetests ist immer von \textbf{weniger} oder mit \textbf{mehr} auszugehen.
\end{bemerkung}

Das \textbf{weniger} oder \textbf{mehr} kann uns leicht den Annahmebereich und den Ablehnungsbereich der Nullhypothese angeben.

\begin{gesetz}{Ablehnungsbereich}{}
  Der Ablehnungsbereich darf maximal mit einer Wahrscheinlchkeit des Signifikanzniveaus eintreten. Dabei wird in der Vorgabe jeweils zwischen $<$ und $\le$ bereits unterschieden.
\end{gesetz}

\begin{gesetz}{Annahmebereich}{}
  Der Annahmebereich ist das Gegenereignis zum Ablehnugsbereich.
\end{gesetz}


\begin{definition}{Einseitiger Hypothesentest}{}
Unter einem \textbf{einseitigen} Hypothesentest verstehen wir einen Ablehnungsbereich von weniger als $\alpha$ auf der einen Seite unserer Binomialverteilung. Beim \textbf{linksseitigen} Hypothesentest liegt der Ablehnungsbereich unserer Nullhypothese im Intervall $[0; k]$, während beim \textbf{rechtsseitigen} Hypothesentest den Ablehnungsbereich im Intervall $[k+1;n]$ liegt.
  \end{definition}


\newpage
\subsubsection{Referenzaufgabe Medikament}
Unser Medikament soll in 80\,\% der Fälle Linderung versprechen. Wir haben Grund zur Annahme, dass es weniger sind.

a) Geben Sie die Nullhypothese und die Alternativhypothese an.

\TNT{4}{Nullhypothese: Das Medikament wirkt in 80\,\% der Fälle.

  Alternativhypothese: Das Medikament wirkt in \textbf{weniger} als 80\,\% der Fälle.
}

b) Begründen Sie, warum Sie hier einen links- bzw. einen rechtsseitigen Hypothesentest anstreben.

\TNT{2.8}{Aufgund der \textbf{weniger}-Alternativhypothese, streben wir einen linksseitigen Hypothesentest an.}

c) Geben Sie eine Entscheidungsregel an ($k$, Ablehnungsbereich, Annahmebereich) für ein Signifikanzniveau von $\alpha \le 1\%$.

\TNT{4}{$P(X \le 69) \approx 0.61\,\%$ und $P(X \le 70) \approx 1.1\,\%$. Somit ist unser $k=69$ für den Ablehnungsbereich $[0;k] = [0;69]$. Der Annahmebereich unserer Nullhypothese liegt bei $[70;100]$ «Genesungen».}


\newpage
\subsubsection{Referenzaufgabe: Rechtsseitiger Hypothesentest}
... coming soon ...
\newpage
